\documentclass[11pt]{article}

\usepackage[margin=1in]{geometry}
\usepackage{amsmath,amssymb,amsthm,mathtools}
\usepackage{graphicx}
\usepackage[hidelinks]{hyperref}

\newtheorem{theorem}{Theorem}
\newtheorem{lemma}{Lemma}
\newtheorem{remark}{Remark}

\newcommand{\R}{\mathbb R}
\newcommand{\Hh}{\mathbb H}
\newcommand{\Sp}{\mathbb S^{N-1}_+}
\newcommand{\Tr}{\operatorname{Tr}}
\newcommand{\norm}[1]{\left\lVert #1\right\rVert}

\title{A Sobolev Trace-Conjunction\\Bourgain--Brezis--Mironescu Formula}
\author{Automated Math Discovery Candidate Proof}
\date{July 23, 2026}

\begin{document}
\maketitle

\begin{center}
\fbox{\parbox{0.91\textwidth}{\textbf{Status.}
Candidate full solution, likely valid pending expert review.
For every \(1\leq p<\infty\), the source formula extends from compactly
supported \(C^1\) functions to the nonhomogeneous Sobolev space
\(W^{2,p}(\R^N_+)\).  No optimality of this class is claimed.}}
\end{center}

\section*{Source and Open Problem}

Let
\[
 \Hh=\R^N_+=\{(x',t):x'\in\R^{N-1},\ t>0\},
 \qquad
 \partial\Hh=\R^{N-1}\times\{0\}.
\]
Van Schaftingen~\cite{VanSchaftingen} proves that if
\(u\in C_c^1(\overline{\Hh})\), then
\begin{equation}\label{eq:source}
\begin{split}
 \lim_{s\to1^-}(1-s)
 \int_{\R^{N-1}}\int_{\Hh}
 \frac{|u(x',0)-u(y)|^p}
 {|(x',0)-y|^{N+sp}}\,dy\,dx'
 &=
 c_{N,p}\int_{\R^{N-1}}|\nabla u(x',0)|^p\,dx',
\\
 c_{N,p}
 &:=
 \frac{\pi^{(N-1)/2}\Gamma((p+1)/2)}
 {p\,\Gamma((N+p)/2)}.
\end{split}
\end{equation}
Open Problem 1 on page 3 asks:
\begin{quote}
\emph{Does (1.8) still hold when \(u\) belongs to some suitable Sobolev
space?}
\end{quote}

\begin{center}
\includegraphics[width=\textwidth]{figures/open_problem_crop.png}
\end{center}

For \(v\in W^{1,p}(\Hh)\), \(1\leq p<\infty\), we use the standard
half-space trace \(\Tr v\in L^p(\R^{N-1})\), characterized by
\[
                  v(\,\cdot\,,t)\longrightarrow\Tr v
                  \quad\text{in }L^p(\R^{N-1})
                  \quad(t\downarrow0).
\]
This characterization, including the endpoint \(p=1\), follows directly
from absolute continuity on normal lines; see also
\cite[Chapter 18]{Leoni}.  If \(u\in W^{2,p}(\Hh)\), both \(\Tr u\) and
\(\Tr\nabla u\) are therefore well defined in \(L^p\).

\section*{Proof Intuition}

Write an interior point relative to a boundary point as
\(y=(x',0)+r\omega\), where \(r>0\) and
\(\omega\in\Sp:=\{\omega\in\mathbb S^{N-1}:\omega_N>0\}\).
The factor \(r^p\) in the first difference converts the energy into
\[
 (1-s)\int_{\Sp}\int_0^\infty
 r^{p(1-s)-1}\norm{Q_r^\omega}_{L^p}^p\,dr\,d\omega,
 \qquad
 Q_r^\omega=
 \frac{u(\,\cdot+r\omega',r\omega_N)-\Tr u}{r}.
\]
The measure \(p(1-s)r^{p(1-s)-1}\,dr\) on \((0,1)\) is a probability
measure concentrating at \(r=0\).

The extra derivative in \(W^{2,p}\) gives exactly the boundary regularity
needed to identify what this Mellin average sees.  Indeed,
\(\nabla u\in W^{1,p}\), so its horizontal slices converge strongly to
\(G=\Tr\nabla u\).  Integrating \(\nabla u\) along the short ray and using
translation continuity gives, uniformly in \(\omega\),
\[
                       Q_r^\omega\longrightarrow G\cdot\omega
                       \quad\text{in }L^p.
\]
The Mellin concentration now yields one factor \(1/p\), while the
hemisphere beta integral supplies the remaining gamma-function factor.
Finally, the nonhomogeneous \(L^p\) norms of \(u\) and \(\Tr u\) control
all radii \(r\geq1\), making that contribution \(O(1-s)\).

\section*{Main Result}

\begin{theorem}\label{thm:main}
Let \(N\geq2\), \(1\leq p<\infty\), and let
\(u\in W^{2,p}(\Hh)\) be real valued.  Then
\begin{equation}\label{eq:main}
\begin{split}
 &\lim_{s\to1^-}(1-s)
 \int_{\R^{N-1}}\int_{\Hh}
 \frac{|\Tr u(x')-u(y)|^p}
 {|(x',0)-y|^{N+sp}}\,dy\,dx'
\\
 &\hspace{25mm}=
 \frac{\pi^{(N-1)/2}\Gamma((p+1)/2)}
 {p\,\Gamma((N+p)/2)}
 \int_{\R^{N-1}}|\Tr\nabla u(x')|^p\,dx'.
\end{split}
\end{equation}
Thus formula \emph{(1.8)} of the source paper holds on the suitable
Sobolev space \(W^{2,p}(\Hh)\), affirmatively answering Open Problem~1.
\end{theorem}

\section*{Uniform Boundary Expansion}

\begin{lemma}\label{lem:ray}
Put \(G=\Tr\nabla u\).  For \(r>0\) and \(\omega\in\Sp\), define
\[
 Q_r^\omega(x')
 =\frac{u(x'+r\omega',r\omega_N)-\Tr u(x')}{r}.
\]
Then
\[
 \lim_{r\downarrow0}\ 
 \sup_{\omega\in\Sp}
 \norm{Q_r^\omega-G\cdot\omega}_{L^p(\R^{N-1})}=0.
\]
In particular,
\(\sup_{0<r<r_0,\,\omega\in\Sp}\norm{Q_r^\omega}_p<\infty\)
for some \(r_0>0\).
\end{lemma}

\begin{proof}
Set \(F=\nabla u\in W^{1,p}(\Hh;\R^N)\).  Strong convergence of horizontal
slices to the trace gives
\[
 \eta(r):=\sup_{0<t\leq r}\norm{F(\,\cdot\,,t)-G}_p
 \longrightarrow0.
\]
Continuity of translations in \(L^p(\R^{N-1})\) gives
\[
 \tau(r):=\sup_{|z|\leq r}\norm{G(\,\cdot+z)-G}_p
 \longrightarrow0.
\]

The \(L^p\)-valued fundamental theorem of calculus along an oblique ray
gives, for almost every \(\omega\in\Sp\),
\begin{equation}\label{eq:rayftc}
 Q_r^\omega
 =\int_0^1
 F(\,\cdot+tr\omega',tr\omega_N)\cdot\omega\,dt
 \quad\text{in }L^p(\R^{N-1}).
\end{equation}
For completeness, one may first integrate from \(\varepsilon\) to \(r\)
for a smooth approximation to \(u\), pass to \(W^{1,p}\), and then let
\(\varepsilon\downarrow0\).  The initial slice converges to \(\Tr u\);
absolute continuity of the \(L^p\)-integral of \(\nabla u\) justifies the
last limit.  This argument is valid for \(p=1\) as well.

For \(0<t\leq1\), invariance of the \(L^p\)-norm under horizontal
translations yields
\[
\begin{split}
 &\norm{
 F(\,\cdot+tr\omega',tr\omega_N)-G}_p
\\
 &\qquad\leq
 \norm{F(\,\cdot,tr\omega_N)-G}_p
 +\norm{G(\,\cdot+tr\omega')-G}_p
 \leq\eta(r)+\tau(r).
\end{split}
\]
Minkowski's integral inequality in \eqref{eq:rayftc} therefore gives
\[
 \norm{Q_r^\omega-G\cdot\omega}_p
 \leq\eta(r)+\tau(r),
\]
independently of \(\omega\).  This proves both assertions.
\end{proof}

\begin{lemma}[Mellin concentration]\label{lem:mellin}
If \(A:(0,1)\to[0,\infty)\) is bounded and
\(A(r)\to A_0\) as \(r\downarrow0\), then
\[
 \lim_{s\to1^-}(1-s)\int_0^1
 r^{p(1-s)-1}A(r)\,dr=\frac{A_0}{p}.
\]
\end{lemma}

\begin{proof}
The measures
\[
        p(1-s)r^{p(1-s)-1}\mathbf1_{(0,1)}(r)\,dr
\]
have total mass \(1\).  For each fixed \(\delta\in(0,1)\), their mass on
\([\delta,1]\) is \(1-\delta^{p(1-s)}\), which tends to zero.  Splitting
the integral at \(\delta\), and then using \(A(r)\to A_0\), proves the
claim.
\end{proof}

\section*{Proof of Theorem \ref{thm:main}}

\begin{proof}
Make the change of variables
\[
                y=(x',0)+r\omega,
                \qquad r>0,\quad\omega\in\Sp.
\]
Polar coordinates and Tonelli's theorem give, for \(0<s<1\),
\begin{equation}\label{eq:polar}
\begin{split}
 E_s(u)
 &:=(1-s)
 \int_{\R^{N-1}}\int_{\Hh}
 \frac{|\Tr u(x')-u(y)|^p}
 {|(x',0)-y|^{N+sp}}\,dy\,dx'
\\
 &=(1-s)\int_{\Sp}\int_0^\infty
 r^{p(1-s)-1}\norm{Q_r^\omega}_p^p\,dr\,d\omega.
\end{split}
\end{equation}

Define
\[
                     A(r)=\int_{\Sp}\norm{Q_r^\omega}_p^p\,d\omega.
\]
Lemma~\ref{lem:ray}, the reverse triangle inequality, and continuity of
\(z\mapsto z^p\) show that \(A\) is bounded near zero and
\begin{equation}\label{eq:A0}
 A(r)\longrightarrow
 A_0:=\int_{\Sp}\norm{G\cdot\omega}_p^p\,d\omega.
\end{equation}
In fact \(A\) is bounded on all of \((0,1)\).  Lemma~\ref{lem:ray}
handles small \(r\); away from zero, use
\[
 \norm{Q_r^\omega}_p
 \leq r^{-1}\bigl(\norm{\Tr u}_p+
 \norm{u(\,\cdot,r\omega_N)}_p\bigr)
\]
and the standard uniform slice bound
\(\sup_{0<t\leq1}\norm{u(\,\cdot,t)}_p<\infty\).
Lemma~\ref{lem:mellin} therefore applies directly, and the part of
\eqref{eq:polar} with \(0<r<1\) converges to \(A_0/p\).

We next compute \(A_0\).  Since
\(\omega\mapsto|e\cdot\omega|^p\) is even, its integral over the upper
hemisphere is one half of its integral over the whole sphere, regardless
of the unit vector \(e\).  The beta integral gives
\begin{equation}\label{eq:angular}
 \int_{\Sp}|a\cdot\omega|^p\,d\omega
 =
 \frac{\pi^{(N-1)/2}\Gamma((p+1)/2)}
 {\Gamma((N+p)/2)}\,|a|^p
 \qquad(a\in\R^N).
\end{equation}
Indeed the left side for \(|a|=1\) is
\[
 \frac12|\mathbb S^{N-2}|
 B\!\left(\frac{p+1}{2},\frac{N-1}{2}\right).
\]
By Fubini's theorem, \eqref{eq:A0} and \eqref{eq:angular} yield
\[
 \frac{A_0}{p}
 =
 \frac{\pi^{(N-1)/2}\Gamma((p+1)/2)}
 {p\,\Gamma((N+p)/2)}
 \int_{\R^{N-1}}|G(x')|^p\,dx'.
\]

It remains to discard the part with \(r\geq1\).  Returning to
\(h=r\omega\), the elementary inequality
\(|a-b|^p\leq2^{p-1}(|a|^p+|b|^p)\) gives
\begin{equation}\label{eq:tail}
\begin{split}
 &\int_{\substack{h_N>0\\|h|\geq1}}
 |h|^{-N-sp}
 \int_{\R^{N-1}}
 |\Tr u(x')-u(x'+h',h_N)|^p\,dx'\,dh
\\
 &\quad\leq 2^{p-1}\norm{\Tr u}_p^p
 \int_{\substack{h_N>0\\|h|\geq1}}|h|^{-N-sp}\,dh
\\
 &\qquad\quad+
 2^{p-1}
 \int_0^\infty \norm{u(\,\cdot,t)}_p^p
 \left[
 \int_{\substack{z\in\R^{N-1}\\|(z,t)|\geq1}}
 |(z,t)|^{-N-sp}\,dz
 \right]dt.
\end{split}
\end{equation}
For \(s\in[1/2,1)\), the first kernel integral is uniformly bounded.
The bracketed integral is also bounded uniformly in \(t>0\) and such
\(s\).  To see this, use
\(|(z,t)|^{-N-sp}\leq|(z,t)|^{-N-p/2}\) on \(|(z,t)|\geq1\).
If \(t\leq1/2\), the restriction forces
\(|z|\geq\sqrt3/2\); if \(t>1/2\), integrate the latter kernel over all
\(z\) and scale by \(t\).  Both bounds are uniform.
Thus the right side of \eqref{eq:tail} is at most
\[
                    C_{N,p}\bigl(
                    \norm{\Tr u}_p^p+\norm{u}_{L^p(\Hh)}^p\bigr).
\]
After multiplication by \(1-s\), the \(r\geq1\) contribution tends to
zero.  Combining this with the small-radius limit proves
\eqref{eq:main}.
\end{proof}

\section*{Verification Notes}

The accompanying script \texttt{code/verify\_constants.py} performs two
independent checks:
\begin{enumerate}
\item it compares the gamma-function expression in
\eqref{eq:angular} with direct one-dimensional beta quadrature for
several noninteger values of \(p\) and several dimensions;
\item it verifies the Mellin normalization on the test family
\(A(r)=A_0+r^q\), for which the integral is exactly
\[
 \frac{A_0}{p}+\frac{1-s}{p(1-s)+q}.
\]
\end{enumerate}
These checks guard against a missing hemisphere factor or a missing
factor \(p\); they are not used as proof.

The principal points for expert review are the \(L^p\)-valued ray identity
\eqref{eq:rayftc}, its uniformity as \(\omega_N\downarrow0\), and the
uniform fixed-height kernel estimate in \eqref{eq:tail}.  Each is written
so that the endpoint \(p=1\) remains visible.

\section*{Novelty Check, Scope, and Limitations}

On July 23, 2026, bounded searches were run in arXiv and general web
indexes for the exact source title, arXiv id \(2502.18140\), the quoted
phrase ``suitable Sobolev space,'' and close combinations of ``trace
conjunction,'' ``Bourgain--Brezis--Mironescu,'' and \(W^{2,p}\).
The searches found the source preprint and its published institutional
repository record, but no later paper answering Open Problem~1 or proving
the extension above.  The run's cheap indexes also had no prior result or
attempt for this arXiv id before reservation.  Novelty confidence is
moderate pending a specialist citation search.

The theorem is a full answer to Open Problem~1 in its literal existence
form: it supplies a suitable Sobolev space for every exponent covered by
the smooth theorem.  It does not claim that \(W^{2,p}\) is sharp.  A larger
space characterized directly by an \(L^p\) boundary trace of the full
gradient would also support the same proof.  The result does not address
the source paper's separate converse problem concerning recovery of a
weak normal derivative from finite limiting conjunction energy.

\section*{Human-Review Recommendation}

Recommend expert verification as a candidate full solution.  In
particular, check:
\begin{enumerate}
\item the trace and oblique-ray fundamental theorem at \(p=1\);
\item the uniform estimate in Lemma~\ref{lem:ray};
\item the far-field estimate \eqref{eq:tail};
\item that the intended reading of ``some suitable Sobolev space'' accepts
the concrete, non-sharp class \(W^{2,p}\).
\end{enumerate}

\begin{thebibliography}{9}

\bibitem{VanSchaftingen}
J.~Van Schaftingen,
\emph{Trace conjunction inequalities},
Academy of Romanian Scientists. Annals. Series on Mathematics and its
Applications \textbf{17} (2025), no.~1, 321--342,
\href{https://arxiv.org/abs/2502.18140}{arXiv:2502.18140},
\href{https://doi.org/10.56082/annalsarscimath.2025.1.321}
{doi:10.56082/annalsarscimath.2025.1.321}.

\bibitem{Leoni}
G.~Leoni,
\emph{A First Course in Sobolev Spaces}, second edition,
Graduate Studies in Mathematics 181, American Mathematical Society,
Providence, RI, 2017.

\end{thebibliography}

\end{document}
